Write $s_+(p)=\max\{0,(n-1)/2-n/p\}$. For $2\le p<\infty$, choose $s_+(p)<\sigma<\alpha$. Choose $\eta\in C_c^\infty((1/2,2))$ with $\sum_{j\in\mathbb Z}\eta(2^jv)=1$ for $v>0$. Apply it to $v=1-\rho$ for $0\le\rho<1$. For $\delta=2^{-j}\le1/8$, the corresponding piece of $(1-\rho^2)_+^\alpha$ equals
\[
\delta^\alpha\psi_\delta((1-\rho)/\delta),\qquad
\psi_\delta(u)=\eta(u)u^\alpha(2-\delta u)^\alpha,
\]
with the expression defined as zero off the support of $\eta$. The profiles have uniform smooth seminorms. The remainder vanishes near $\rho=1$ and equals $(1-\rho^2)^\alpha$ near zero, so its radial extension is smooth and compactly supported. The kernel of this remainder has an integrable radial decreasing majorant. The supremum of its dilations is therefore bounded pointwise by a constant times the Hardy--Littlewood maximal function. Apply \ref{ms:lemma-10-2} to the remaining pieces. The sum of their maximal $L^p$ norms is bounded by
\[
C\sum_{j\ge3}2^{-j(\alpha-\sigma)}\|f\|_p.
\]
This proves the maximal bound on Schwartz functions, and gives the fixed-radius bound there since $\delta(p)=s_+(p)$ in this range.

Each fixed $B_r^\alpha$ is self-adjoint on Schwartz functions, and $\delta(p)=\delta(p')$ for conjugate exponents. Duality proves the fixed-radius bound for $1<p<2$. At $p=1,\infty$ the hypothesis is $\alpha>(n-1)/2$. Radial Fourier inversion gives the unit-radius kernel as a constant times $|x|^{-n/2-\alpha}J_{n/2+\alpha}(|x|)$, where $J_\nu$ is the Bessel function of the first kind. The kernel is bounded at zero and is $O(|x|^{-(n+1)/2-\alpha})$ at infinity, so its dilates $r^nK(rx)$ have a common $L^1$ norm. Density for finite $p$ and convolution at $p=\infty$ give the stated extensions.

For Schwartz functions with compact Fourier support, the symbols converge to one in all derivatives on the support as $r\to\infty$. The means then converge in the Schwartz topology. Density and the fixed-radius bound give norm convergence for $p<\infty$. For the almost everywhere and maximal claims, choose such $f_j$ with $\sum_j\|f-f_j\|_p<\infty$. The maximal bound on Schwartz functions implies that $\sum_j\sup_{r>0}|B_r^\alpha(f_{j+1}-f_j)|$ is finite almost everywhere. At every such point $B_r^\alpha f_j$ converges uniformly in $r>0$ to a continuous function of $r$, agreeing for each fixed $r$ almost everywhere with the bounded $L^p$ extension. Its supremum over $r$ is measurable by continuity, and Fatou's lemma gives the maximal estimate for $f$. Finally $f_j(x)\to f(x)$ almost everywhere, and uniform approximation in $r$ transfers convergence as $r\to\infty$ from $f_j$ to $f$.
